\chapter{Artifact and Reproducibility Description}
\label{app:artifacts}

All software, configurations, and derived data of this thesis are archived
in a single repository (\code{github.com/mij001/cuvslam-stack}, MIT
license). The artifact is organized so that \emph{every table in this
document regenerates from committed CSV data using only the Python standard
library no GPU, no datasets, no vendor tools required}; capture (which
does require the hardware) is separated from analysis by construction.

\section{Layout}

\begin{description}
\item[\code{configs/}] The configuration regime: a small set of human-owned
base TOML files plus a deterministic mutation pass that materializes the
full experiment matrices (141 accuracy configurations, 192
profiling-coverage variants). One TOML fully describes a run: dataset,
camera rig, pipeline mode, feature toggles, and the ground-truth evaluation
block.
\item[\code{profiling/harness/}] The capture entry point: one command runs
any configuration under nsys, ncu (curated metric set, bounded windows), or
NVBit (windowed, kernel-filtered address traces), writing versioned result
directories with full provenance (clock locks, driver/tool versions,
command lines). A workload-\emph{adapter} interface makes the harness
generic: cuVSLAM is one adapter; a command adapter runs arbitrary GPU
programs through the identical pipeline.
\item[\code{profiling/analysis/}] The stdlib-only analysis library:
screening, classification (the decision tree of \S\ref{sec:met-tree}, with
thresholds defined once and imported live into documentation), roofline,
locality (reuse-distance CDFs with memory-space filtering), attribution
(the journal-trace join), campaign synthesis, substrate verdicts, residual
ranking, and the DAMOV-parity checks. The GPU-free test suite fabricates
known inputs for every module.
\item[\code{profiling/calibration/}] The known-truth archetype suite
(\S\ref{sec:char-calibration}): eight kernels, one per class, with the
blind-classification driver.
\item[\code{profiling/validation/}] The clock-domain intervention scripts
(\S\ref{sec:met-parity}) and the profiler-neutrality drivers.
\item[\code{patches/}] The instrumentation as reviewable diffs: the
TaggedAllocator + NVTX enablement patch against the cuVSLAM source tree,
and the NVBit \code{mem\_trace} extensions (launch-window and kernel-name
gating; the allocation-event sidecar).
\item[\code{reports/}, \code{profiling/reports/}] Dated, immutable result
directories: the 27-sequence characterization campaign, the locality study
and its space-filtered re-derivation (both the superseded and corrected
tables are retained, with the correction marked in place), the attribution
campaign, the 141-run accuracy matrix, the 192-variant neutrality campaign,
the substrate synthesis, the energy and host-I/O measurements, and the
DAMOV-validation artifacts (cross-device agreement, hierarchical-clustering
agreement).
\item[\code{dashboard/}] The evidence explorer: an interactive application
presenting every finding with its computed number, the per-run reasoning
chain (time share $\rightarrow$ metrics vs.\ thresholds $\rightarrow$ the
decision-tree rule that fired $\rightarrow$ verdict), and a methodology tab
whose formulas and thresholds are stamped from the analysis source at build
time so interface and code cannot disagree.
\end{description}

\section{Environment}

Primary testing hardware: NVIDIA RTX~2000 Ada (\code{sm\_89}), driver
575.64.05, CUDA 12.9.1, Nsight Compute 2025.2.1.3, Nsight Systems 2025.3.2,
and NVBit 1.8. Maintaining this exact combination of software and driver
versions is critical, since the binary instrumentation specifies the driver
version, which, in its turn, determines the profiler versions. To ensure
compatibility, an environment doctor checks the target system and suggests
a one-line fix for any incompatibilities detected. GPU clocks are locked
(1620/7001\,MHz) for every measurement; ceilings are re-measured per
instrument.

\section{Reproduction paths}

\textbf{Analysis-only (no GPU):} clone; run the test suite; regenerate any
table in this thesis from the committed CSVs via the corresponding
\code{analysis} module. \textbf{Full capture:} configure a target with the
environment doctor; run the regime's campaign drivers (all are resumable
and idempotent re-running skips completed steps); the coverage audit
then proves the captures complete before analysis will consume them.
\textbf{Instrumented builds:} apply the committed patch to the cuVSLAM
source tree and build with the pinned container recipe; the baseline and
instrumented binaries are kept side by side, and the QoR check of
\S\ref{sec:val-neutrality} re-verifies output equivalence on any rebuild.
